% Korean edition of sections/s2_2_auth.tex.
\seclabel
\claim{\secnum{2.2}{sec:auth}\quad 액세스 토큰은 900\,s 유효하고, 갱신도 로그아웃도 7일짜리 쿠키 묶음을 바꾸지 않는다}

\lead{로그인하면 응답 하나에 액세스 토큰과 포트폴리오 목록이 함께 온다.} \code{{"email", "password"}}로 \code{POST /api/Auth/Login}을 보내면 HTTP 200과 함께 JWT 형식의 435자 액세스 토큰(\code{accessToken}), 포트폴리오 6개(\code{pfs}), 대회 16개(\code{ctsts}), 서버의 현재 영업일(\code{currBizDay})이 온다. 같은 응답에 개인 정보(연락처) 필드도 들어 있으므로, 래퍼는 필요한 필드만 남기고 이 응답 전체를 로그에 남기지 않는다.

\lead{로그인 응답은 유효기간이 서로 다른 자격 증명 5개를 설정하며, 래퍼에는 액세스 토큰과 쿠키 묶음만 필요하다(\figref{creds}).} 액세스 토큰은 모든 REST 요청과 피드 협상의 \code{Authorization: Bearer} 헤더에 들어간다. \code{refreshToken} 쿠키는 루트 경로의 다른 쿠키와 함께 보내 새 액세스 토큰을 받는 데 쓴다.

% Korean edition of figs/creds.tex.
\fig[wide]{creds}{액세스 토큰은 900\,s 유효하고, 래퍼는 쿠키 묶음 전체를 저장한다}[2026년 11월 24일 05:38:54 KST 로그인이 돌려준 자격 증명과 래퍼가 각각을 저장하는 곳]{%
{\tabfont\begin{tabularx}{\linewidth}{>{\RaggedRight\arraybackslash}p{24mm} >{\RaggedRight\arraybackslash}p{20mm} >{\RaggedRight\arraybackslash}p{50mm} >{\RaggedRight\arraybackslash}p{20mm} >{\RaggedRight\arraybackslash}X}
\thd{이름} & \thd{전달 위치} & \thd{유효기간} & \thd{저장 위치} & \thd{용도}\\
\hl\code{accessToken} & 응답 본문 & \textbf{900\,s}, JWT \code{exp} $-$ \code{iat}. 클레임: \code{Id}, \code{exp}, \code{iat}, \code{nbf}, \code{role} & 메모리만 & 모든 REST 요청의 \code{Authorization: Bearer} 헤더\\\rowsep
\hl\code{refreshToken} & 쿠키 & \textbf{7일}, 12월 1일 05:38까지. 경로 \code{/} & \code{session.json} & \code{Auth/SilentRefresh}로 갱신, 루트 경로의 다른 쿠키와 함께 보냄\\\rowsep
\code{asset_token} & 쿠키 & 1시간. 경로 \code{/assets/app/} & \code{session.json} & 로그인한 사용자에게만 내려오는 페이지 코드 다운로드\\\rowsep
\code{userid} & 쿠키 & 7일 & \code{session.json} & 페이지 JavaScript가 읽는 사용자 ID\\\rowsep
\code{nxdid} & 쿠키 & 1년 & \code{session.json} & 기기 식별자로 추정\\
\end{tabularx}}}{음영 행은 유효기간이 래퍼의 일정을 정하는 자격 증명 2개다. 액세스 토큰 발급 후 720\,s(서버 시계)가 지난 뒤 첫 REST 요청이 토큰을 갱신하고, \code{refreshToken}의 남은 기간이 1일 미만이면 래퍼가 다시 로그인한다. \code{refreshToken}만으로 갱신하면 401로 거부되므로, 래퍼는 로그인과 갱신 때마다 쿠키 묶음의 모든 쿠키를 \code{~/.config/pytimefolio/session.json}(권한 0600)에 기록한다. 액세스 토큰은 디스크에 쓰지 않는다. 쿠키 값은 기록하지 않았고 이름, 경로, 유효기간만 기록했다. \code{nxdid}의 용도는 이름과 유효기간으로 추정했다.}{서버 응답 헤더; \code{pytimefolio/auth.py}; DART Quant Team}


\lead{갱신에는 쿠키 묶음 전체가 필요하며, 갱신해도 7일 만료 시점은 그대로다.} 루트 경로 쿠키를 싣고 \code{Authorization} 헤더 없이 \code{POST /api/Auth/SilentRefresh}를 보내면 HTTP 200과 새 액세스 토큰이 온다. \code{refreshToken}만 보내면 401로 거부되므로 래퍼는 쿠키 묶음 전체를 보관하고 보낸다(D4). \code{refreshToken}의 만료 시각은 호출 전후 모두 11월 30일 20:38:54 UTC였고, 응답은 새 \code{refreshToken}을 설정하지 않았다. 즉 갱신은 만료를 뒤로 밀지도, 토큰을 교체하지도 않으며, 같은 쿠키로 반복한 갱신은 모두 성공했다. 장중에 같은 쿠키로 \code{SilentRefresh} 2건을 동시에 보내자 둘 다 200을 받았고, 바로 뒤에 보낸 세 번째도 200을 받았으며, 어느 것도 새 \code{refreshToken}을 설정하지 않았다. 동시 갱신은 서로를 무효화하지 않는다.

\lead{서버는 만료된 액세스 토큰을 300\,s 동안 받아준다.} 첫 테스트에서 거부 시점은 만료 후 3.5\textasciitilde 5.6분 사이로 좁혀졌다. 서버 시계로 08:16:43에 발급된 두 번째 토큰을 만료 후 5분 직전과 직후에 다시 보냈다. 서버의 \code{Date} 헤더 기준으로 만료 후 4분 57초에는 받아들여졌고 5분 07초에는 거부됐다(\figref{token}). 거부 시점은 이 10\,s 구간 안에 있으며, ASP.NET JWT 검증의 기본 시계 오차 허용치 300\,s[2]와 맞는다. 401의 문구도 Microsoft의 .NET 인증 라이브러리에서 나온다. \figref{token}의 두 번째 줄, 즉 토큰보다 오래 살아남은 열린 WebSocket은 실시간 피드에 해당한다(\ref{sec:feed}절).

% Korean edition of figs/token.tex.
\fig[wide]{token}{서버는 만료 후 300\,s 동안 액세스 토큰을 받아주며, 열린 WebSocket은 토큰보다 오래 산다}[2026년 11월 24일 세 테스트를 각 액세스 토큰의 만료 시점(x = 0)에 맞춰 정렬: 서버 시계로 08:16:43에 발급된 토큰의 REST 요청, 07:21:41 KST에 만료된 토큰의 WebSocket 하나, 05:38:54 로그인에서 설정된 갱신 쿠키]{%
\begin{tikzpicture}[x=6mm,y=1mm,lab/.style={font=\fontsize{7}{8.4}\selectfont},lane/.style={font=\fontsize{7.5}{9}\selectfont\bfseries,anchor=east,text width=24mm,align=right}]
% ---- lane labels
\node[lane] at (-15.25,0) {REST 요청};
\node[lane] at (-15.25,-17) {열린 WebSocket};
\node[lane] at (-15.25,-31) {갱신 쿠키};
% ---- expiry line
\draw[navy,dashed,line width=0.5pt] (0,-36) -- (0,15.5);
\node[lab,anchor=south east] at (-0.1,12.2) {만료};
% ---- lane 1: REST
\fill[navy] (-15,-2.5) rectangle (0,2.5);
\fill[navyL] (0,-2.5) rectangle (4.95,2.5);
\fill[pattern=north east lines,pattern color=navyM] (4.95,-2.5) rectangle (5.1167,2.5);
\draw[navyM,line width=0.3pt] (4.95,-2.5) rectangle (5.1167,2.5);
\fill[navyXL] (5.1167,-2.5) rectangle (9,2.5);
\node[lab,text=white] at (-7.5,0) {유효: 발급 후 900\,s};
\node[lab] at (2.475,0) {허용};
\node[lab,text=greyD] at (7.06,0) {거부};
\node[lab,anchor=south west,text=greyD] at (-15,2.8) {발급};
% 200 and 401 marks
\fill[navy] (4.95,6) circle (0.75mm); \draw[greyM,dashed,line width=0.3pt] (4.95,5.2) -- (4.95,2.5);
\node[lab,anchor=east] at (4.75,6) {+4:57에 200};
\draw[navy,line width=0.8pt] ($(5.1167,11)+(-0.8mm,-0.8mm)$) -- ($(5.1167,11)+(0.8mm,0.8mm)$) ($(5.1167,11)+(-0.8mm,0.8mm)$) -- ($(5.1167,11)+(0.8mm,-0.8mm)$);
\draw[greyM,dashed,line width=0.3pt] (5.1167,10) -- (5.1167,2.5);
\node[lab,anchor=west] at (5.35,11) {+5:07에 401};
% tolerance bracket (timeout style)
\draw[greyD,line width=0.4pt] (0,-4) -- (0,-5) -- (5,-5) -- (5,-4);
\node[lab,anchor=north,text=greyD] at (2.5,-5.2) {허용치 300\,s};
% ---- lane 2: WebSocket replies, one tick per reply
\draw[greyM,line width=0.3pt] (-11.42,-17) -- (9,-17);
\foreach \i in {0,...,381}{\draw[navy,line width=0.25pt] ({-11.42+\i*0.05052},-19) -- ({-11.42+\i*0.05052},-15);}
\node[lab,anchor=south west] at (-11.42,-14.6) {응답마다 눈금 1개, 3\,s마다 폴링};
\node[lab,anchor=south east] at (7.83,-14.6) {끝까지 응답 382개};
\draw[navy,line width=0.9pt] ($(8.3167,-17)+(-1mm,-1mm)$) -- ($(8.3167,-17)+(1mm,1mm)$) ($(8.3167,-17)+(-1mm,1mm)$) -- ($(8.3167,-17)+(1mm,-1mm)$);
\node[lab,anchor=north east,align=right] at (9.2,-19.6) {07:30(+8:19) 새 협상,\\ 같은 토큰: 401};
% ---- lane 3: refresh cookie
\fill[navyM] (-15,-33) rectangle (9,-29);
\draw[navyM,line width=1pt,-{Stealth[length=2mm]}] (9,-31) -- (10.1,-31);
\node[lab,text=white,anchor=west] at (-14.7,-31) {로그인 때 설정, 7일 뒤인 12월 1일 05:38:54 KST에 만료. 갱신해도 만료는 그대로};
% ---- axis
\draw[greyM,line width=0.4pt] (-15,-38) -- (9,-38);
\foreach \m/\t in {-15/−15,-12/−12,-9/−9,-6/−6,-3/−3,0/0,3/+3,6/+6,9/+9}{\draw[greyM] (\m,-38) -- (\m,-39); \node[lab,text=greyD,anchor=north] at (\m,-39) {\t};}
\node[lab,text=greyD,anchor=north] at (-3,-42.2) {액세스 토큰 만료 후 경과 시간(분)};
\end{tikzpicture}
}{음영: 남색은 액세스 토큰이 유효한 구간, 연한 남색은 만료됐지만 받아주는 구간, 빗금은 REST 거부 시점이 있는 10초 구간, 옅은 색은 REST 요청이 거부되는 구간이다. 허용치 경계는 WebSocket 테스트와 다른 토큰으로 한 두 번째 테스트에서 나왔으며 그 토큰의 만료 시점 기준으로 배치했다. PC 시계가 3\,s 빨랐으므로 시간은 서버의 \code{Date} 헤더로 쟀다. 401에는 \code{www-authenticate: Bearer error="invalid_token", error_description="The token expired at ..."}(UTC)가 실려 있었다. WebSocket 눈금 하나는 같은 연결의 응답 하나이며, 눈금 위치는 테스트 스크립트가 1분마다 남긴 첫 로그 줄과 마지막 로그 줄 사이에 고르게 배치했다. X는 거부된 요청이다. 쿠키 만료는 갱신 전후의 \code{Set-Cookie}에서 읽었다.}{DART Quant Team}


\lead{우리 PC 시계가 서버보다 3\,s 빨라서, 래퍼는 서버 시계로 만료를 계산한다.} 두 번째 테스트에서 서버 \code{Date} 헤더가 23:36:40 UTC일 때 로컬 시계는 08:36:43이었다. 자기 시계를 JWT의 \code{exp}와 비교하는 래퍼는 이 오차를 그대로 물려받는다. 그래서 래퍼는 모든 응답의 \code{Date} 헤더에서 시차를 구하고 최근 9개 시차의 중앙값을 쓴다. 헤더 해상도가 1\,s이므로 표본마다 0.5\,s를 더한다.

\lead{갱신은 필요할 때 일어난다. 토큰 발급 후 720\,s 이상 지나서 나가는 첫 REST 요청이 토큰을 갱신한다.} 타이머는 없다. 래퍼는 REST 요청마다 서버 시각을 토큰의 \code{iat}와 비교하고, 720\,s가 지났으면 먼저 \code{SilentRefresh}를 호출한 뒤 새 토큰으로 요청을 보낸다. 900\,s 만료 전에 남는 180\,s는 느린 갱신을 대비한 여유이며, 일정은 바뀔 수 있는 서버 설정인 300\,s 허용치에 기대지 않는다. 열린 WebSocket은 토큰이 만료돼도 계속 동작하므로(\ref{sec:feed}절) 실시간 피드는 연결하거나 다시 연결할 때만 토큰이 필요하다. 동시 요청은 갱신 하나를 공유한다. 오프라인 테스트에서 토큰을 동시에 5번 요청하자 \code{SilentRefresh}는 1번 호출됐다.

\lead{잘못된 요청과 로그인 실패는 스택의 서로 다른 부분에서 서로 다른 오류를 받는다.} JSON이 깨진 요청 본문(문자열 안의 날 캐리지 리턴)은 nginx 형식의 HTML 400 페이지를 받았다. 올바른 JSON에 틀린 자격 증명을 보내면 API가 401을 돌려준다. 400은 요청을 만든 코드의 결함을 뜻해 다시 보내도 실패하므로, 래퍼는 재시도 없이 예외를 던진다. 401은 갱신이나 새 로그인으로 해결된다.

\lead{401을 받으면 래퍼는 토큰을 한 번 갱신하고 요청을 다시 보낸다. ASP.NET Core 문서의 미들웨어 순서에 근거해, 서버가 요청을 처리하기 전에 거부했다고 가정한다[9].} ASP.NET Core는 요청을 처리하는 핸들러보다 인증과 인가 단계를 먼저 실행하고, 여기서 실패한 요청은 핸들러에 닿지 않고 401을 받는다. 이에 따르면 다시 보낸 요청은 주문이라 해도 두 번 실행될 수 없다. 주문 요청에서 401을 본 적은 없으므로 이는 문서에만 근거한다. 갱신 뒤 두 번째 401은 예외로 던진다.

\lead{다시 로그인해도 기존 로그인 세션은 계속 동작한다.} 같은 계정으로 두 번째 로그인을 한 뒤에도 첫 로그인 세션의 쿠키로 토큰이 갱신됐고, 브라우저는 페이지를 새로 고친 뒤에도 로그인 상태였다. 따라서 래퍼와 웹 페이지를 동시에 쓸 수 있다. 다만 로그인할 때마다 7일 동안 살아 있는 로그인 세션이 하나씩 늘어난다.

\lead{로그아웃은 브라우저 쿠키를 지우지만, 저장해 둔 쿠키 묶음은 7일 만료까지 계속 갱신된다.} \code{POST /api/Auth/Logout}은 본문 \code{{}}와 함께 200을 돌려주고, \code{refreshToken}, \code{userid}, \code{asset_token}을 1970년 만료일로 비우는 \code{Set-Cookie} 헤더를 보냈다. 호출 전에 쿠키 묶음 사본을 저장해 두었다. 로그아웃 뒤에도 그 사본으로 토큰이 갱신됐고(\code{SilentRefresh} 200), 로그아웃 전에 발급된 액세스 토큰도 REST에서 계속 동작했으며(200), 같은 계정의 다른 로그인 세션 2개도 계속 동작했다. 따라서 저장된 쿠키 묶음은 브라우저에서 무슨 일이 있든 7일짜리 자격 증명으로 동작하며, 래퍼는 이를 전제로 설계됐다.

\lead{래퍼는 쿠키 묶음을 저장하고 비밀번호는 절대 저장하지 않는다.} 쿠키 묶음은 로그인과 갱신 때마다 \code{~/.config/}\allowbreak\code{pytimefolio/}\allowbreak\code{session.json}에 권한 0600으로, 원자적으로, 쿠키별 만료 시각과 함께 기록된다. 비밀번호는 호출자가 생성자에 넘기며(예: 환경 변수에서), 새로 로그인할 때만 쓴다. 저장된 쿠키 묶음의 보안 영향은 \ref{sec:design}장에 정리한다.
\verdict{래퍼는 서버 시계로 720\,s에 갱신한다. 401을 받으면 한 번 갱신하고 요청을 다시 보내며, 주문도 마찬가지다. 401은 주문 처리기보다 먼저 오기 때문이다(D3, D36).}
\verdict{래퍼는 갱신이 실패하거나 저장된 쿠키 묶음의 남은 기간이 1일 미만일 때만 호출자의 비밀번호로 로그인하고, 그 밖에는 쿠키 묶음으로 갱신한다(D5, D17). 쿠키 묶음 파일에 쿠키별 만료 시각을 남기므로 다시 시작한 뒤에도 1일 규칙이 적용된다. 쿠키 묶음은 저장소 밖에 소유자 전용 권한으로 두며, 종료할 때 \code{Logout}을 호출하지 않는다. 호출 뒤에도 저장된 쿠키 묶음은 서버에서 계속 갱신되는 반면, 호출은 로컬 사본을 비워 다음 시작 때 새 로그인과 새 로그인 세션을 만들기 때문이다(D16).}
